\documentclass{article}
\usepackage{tcee-repaso}
\tcetema[Objeto y métodos de la economía]{3.A.1}{Objeto y métodos de la ciencia económica. Cuestiones y debates actuales, con especial referencia a la economía conductual.}
\tcefecha{04/11/2026}
\begin{document}
\cabecerarepaso
\begin{columnas}

% ==========================================================================
\bloque{Esquema del cante}
\esquemacante{%
\apartadocante{Introducción: «a toro pasado\ldots»}{2,5}
\apartadocante{I. Objeto de la ciencia económica}{6}
\subcante{1. Definiciones históricas}{2,5}
\subcante{2. Positiva y normativa; imperialismo}{2}
\subcante{3. Ramas y evolución del objeto}{1,5}
\apartadocante{II. Métodos de la ciencia económica}{9}
\subcante{1. ¿Es la economía una ciencia?}{1}
\subcante{2. Método científico y economía}{5,5}
\subcante{3. Debates metodológicos}{2,5}
\apartadocante{III. Debates actuales; economía conductual}{11,5}
\subcante{1. Racionalidad y economía conductual}{7}
\subcante{2. Experimentos}{2}
\subcante{3. Innovaciones metodológicas}{2,5}
\apartadocante{Conclusión}{1}}
\emph{Reparto} de \hyperref[rep:II.2]{II.2}: deducción e inducción \tiempo{1}, \emph{Methodenstreit} \tiempo{0,5}, \autor{Popper} \tiempo{1}, \autor{Kuhn} \tiempo{1,5}, \autor{Lakatos} y \autor{Feyerabend} \tiempo{1}, escuelas \tiempo{0,5}. De \hyperref[rep:III.1]{III.1}: idea y \autor{Simon} \tiempo{1}, $\beta$-$\delta$ \tiempo{2,5}, perspectivas \tiempo{1}, otros sesgos \tiempo{1}, \emph{nudges} y finanzas \tiempo{1}, valoración \tiempo{0,5}. \emph{Pizarra}: refrán en una esquina (vuelve en la idea final). \emph{I}: escalera de definiciones (izda.), tabla positiva/normativa/arte columna a columna (dcha.), árbol micro-macro con la flecha de la microfundamentación. \emph{II}: 5 hitos de la línea del tiempo arriba; ciclo de \autor{Kuhn} y núcleo-cinturón de \autor{Lakatos}, imprescindibles; se borra para \hyperref[rep:II.3]{II.3} (de palabra; 7 cualidades en lista). \emph{III}: función $\beta$-$\delta$ recuadrada arriba a la izda. hasta el final de \hyperref[rep:III.1]{III.1}, con el descuento y la función de valor al lado.

% ==========================================================================
\bloque[2,5]{Introducción}
\begin{itemize}
\item \textbf{Enganche}: «A toro pasado, todos somos Manolete». Nov. 2008: Isabel~II en la LSE: ¿por qué nadie la vio venir? Jul. 2009, British Academy: fallo de la imaginación colectiva de mucha gente brillante. 3 preguntas: ¿qué estudia la economía, cómo sabe lo que sabe y qué hace cuando la realidad la desmiente?
\item \textbf{Relevancia}: \emph{teoría}: objeto y método deciden qué cuenta como conocimiento; cada escuela = un «qué» con un «cómo» \vertema{3.A.2} \vertema{3.A.7}. \emph{Política}: la ciencia social con más peso en lo público; evaluación con experimentos y \emph{nudges} \vertema{3.B.28}. \emph{Actualidad}: 2008 reabre los modelos macro; microdatos y aprendizaje automático; Nobel a \autor{Kahneman} (2002), \autor{Thaler} (2017), \autor{Banerjee}-\autor{Duflo}-\autor{Kremer} (2019), \autor{Card}-\autor{Angrist}-\autor{Imbens} (2021). 2 de los 4 motivos de \autor{Schumpeter} (1954) para la historia del pensamiento.
\item \textbf{Contextualización}: \emph{oikonomía}, administración de la casa (\autor{Jenofonte}, s.~IV a.~C.); siglos dentro de moral, derecho y teología (precio justo, usura). «Economía política»: \autor{Montchrestien} (1615); en inglés, \autor{Steuart} (1767). Se independiza con \autor{Smith} (1776), tras los fisiócratas; \emph{economics}: \autor{Marshall} (1890). Nobel desde 1969. «Economía es lo que hacen los economistas» (atribuida a \autor{Viner}).
\item \textbf{Problemática}: ¿qué estudia y dónde están sus fronteras? ¿Cómo construye y contrasta teorías y cómo cambian? ¿Qué es un buen modelo; sirve con supuestos irreales? ¿Qué cambia si las personas no se comportan como supone la teoría?
\item \textbf{Estructura}: \hyperref[rep:I]{I}, objeto (de la riqueza a la elección; positiva y normativa; ramas); \hyperref[rep:II]{II}, métodos (¿ciencia?; de deducción e inducción a \autor{Popper}, \autor{Kuhn} y \autor{Lakatos}; modelos y supuestos); \hyperref[rep:III]{III}, debates actuales (conductual con $\beta$-$\delta$ como modelo central; experimentos; innovaciones tras 2008).
\end{itemize}

% ==========================================================================
\bloque[6]{I. Objeto de la ciencia económica\label{rep:I}}
\begin{ideaclave}
La definición pasa de la \emph{riqueza} (clásicos) a la \emph{elección} con recursos escasos (\autor{Robbins}) y a un \emph{modo de razonar} aplicable a cualquier conducta (\autor{Becker}). Detrás, el debate: ¿describe lo que es o prescribe lo que debe ser?
\end{ideaclave}
\begin{itemize}
\item \concepto{Objeto}: fenómenos que estudia una ciencia; \concepto{método} (\emph{méthodos}, camino): cómo llega a conocerlos y contrastarlos. Escuela (\autor{Schumpeter}): líder, seguidores, cuerpo doctrinal.
\end{itemize}

\sub{I.1. Definiciones históricas\label{rep:I.1}}
\begin{itemize}
\item \emph{Clásicos: la riqueza}. \autor{Smith} (1776, libro~IV): ciencia del legislador con 2 objetos: ingreso abundante al pueblo y suficiente al Estado; ciencia práctica, más \emph{debe ser}. \autor{Say}, \autor{Senior} (1836): \emph{leyes} de producción, distribución y consumo de la riqueza.
\item \autor{Mill} (1836): leyes de los fenómenos sociales de producir riqueza, \emph{en la medida en que no los modifica otro objeto}. Separa ciencia de arte; deja fuera las leyes técnicas.
\item \emph{Neoclásicos: las personas}. \autor{Marshall} (1890): «estudio de la humanidad en los asuntos ordinarios de la vida»; requisitos materiales del bienestar; germen del individualismo metodológico.
\item \emph{\autor{Robbins} (1932): la elección}: \concepto{ciencia que estudia el comportamiento humano como una relación entre fines y medios escasos que tienen usos alternativos}. 3 elementos: \emph{escasez}, \emph{usos alternativos} (de los medios), \emph{coste de oportunidad}. Definición \emph{formal}: un aspecto de toda conducta $\Rightarrow$ abre el imperialismo. Encaja con optimizar con restricciones \vertema{3.A.8}.
\item \emph{Mecanismos de asignación} (\autor{Landreth}-\autor{Colander}): fuerza, tradición, autoridad, mercado. Cálculo socialista: \autor{Mises} (1920) sin precios no hay cálculo; \autor{Lange} (1936) tanteo; \autor{Hayek} (1945) información dispersa.
\item \emph{\autor{Becker}: modo de razonar}: maximización, preferencias estables, equilibrio, aplicados a educación, familia, discriminación, delito (1976); Nobel 1992. \emph{Freakonomics} (\autor{Levitt}-\autor{Dubner}, 2005).
\item \emph{\autor{Viner}}: «lo que hacen los economistas»: no está en sus obras; la recoge \autor{Boulding} (1941). ¿Quién es economista? \autor{Ferreira} (2013): quien publica con revisión por pares; \autor{Tirole} (2017): como el médico, diagnostica y trata.
\end{itemize}

\sub{I.2. Positiva y normativa; imperialismo económico\label{rep:I.2}}
\begin{itemize}
\item \emph{Triple distinción}: J.~N.~\autor{Keynes} (1891), recuperada por \autor{Colander} (1992). \concepto{Positiva}: lo que \emph{es}, causa-efecto contrastable («al pan, pan»). \concepto{Normativa}: lo que \emph{debe ser}, con juicios de valor (bienestar \vertema{3.A.22}). \concepto{Arte}: «sistema de reglas para alcanzar un fin dado».
\end{itemize}
\begin{center}\footnotesize\setlength{\tabcolsep}{2.5pt}
\begin{tabular}{@{}>{\raggedright\arraybackslash}p{0.15\linewidth}>{\raggedright\arraybackslash}p{0.25\linewidth}>{\raggedright\arraybackslash}p{0.25\linewidth}>{\raggedright\arraybackslash}p{0.29\linewidth}@{}}
\toprule
 & \textbf{Positiva} & \textbf{Normativa} & \textbf{Arte}\\
\midrule
Pregunta & Lo que es & Lo que debe ser & Cómo conseguirlo\\
Valores & Se evitan & Su contenido & Fines dados\\
Método & Abstracto, contrastable & Ético & Integra lo político e institucional\\
Salario mín. & ¿Efecto en el empleo? & ¿Es justo? & ¿Qué nivel sin destruir empleo?\\
\bottomrule
\end{tabular}
\end{center}
\begin{itemize}
\item Arte: escuela histórica, \autor{Marshall} (estela de \autor{Smith}); positiva: ortodoxia moderna (\autor{Ricardo}).
\item \emph{Separación hechos-valores}: \autor{Robbins} (1932): economía = hechos, ética = valoraciones. \autor{Friedman} (1953): positiva independiente de la ética; la normativa se apoya en ella (su regla monetaria) \vertema{3.A.5}.
\item \emph{Crítica}: \autor{Myrdal} (1953): conceptos con valoraciones implícitas (el criterio de \autor{Pareto} ignora el reparto): hacerlas explícitas. \emph{Werturteilsstreit}: \autor{Weber} (\emph{Wertfreiheit}) vs. \autor{Schmoller}.
\item \emph{Lectura conductual}: \autor{Thaler} (1980): la teoría del consumidor es \emph{normativa}; propone una teoría \emph{positiva} de la elección. La reina pedía positiva; lo reprochado después fue arte (regulación financiera).
\item \concepto{Imperialismo económico}: con \autor{Robbins} el objeto no se limita al mercado; \autor{Becker} \vertema{3.A.25}; \autor{Lazear} (2000) lo celebra. Crítica: «quien mucho abarca poco aprieta»: sin instituciones ni psicología. La conductual, movimiento contrario: importa de la psicología.
\end{itemize}

\sub{I.3. Ramas y evolución del objeto\label{rep:I.3}}
\begin{itemize}
\item \emph{Micro}: decisiones y coordinación en mercados; cimientos neoclásicos (\autor{Walras}, \autor{Marshall}) \vertema{3.A.3}; de ahí bienestar, pública, información, trabajo, comercio, finanzas, conductual.
\item \emph{Macro}: la economía en su conjunto; nace con la \emph{Teoría general} (1936) \vertema{3.A.4}: la Gran Depresión cambia el objeto. Frontera difuminada: tras \autor{Lucas}, \concepto{microfundamentación}; diferencia más de preguntas que de método \vertema{3.A.6}; en método debe más a \autor{Fisher} que a \autor{Keynes} \vertema{3.A.29}.
\item \emph{Debates transversales}: papel del Estado \vertema{3.A.23}, neutralidad del dinero \vertema{3.A.5}, existencia y estabilidad del equilibrio \vertema{3.A.21}.
\end{itemize}

% ==========================================================================
\bloque[9]{II. Métodos de la ciencia económica\label{rep:II}}
\begin{ideaclave}
Método hipotético-deductivo contrastable, con 2 problemas propios: experimenta poco y sus «partículas» piensan. La filosofía de la ciencia pasa de buscar un criterio de verdad a explicar cómo cambian las teorías; la economía recorre el mismo camino.
\end{ideaclave}

\sub{II.1. ¿Es la economía una ciencia?\label{rep:II.1}}
\begin{itemize}
\item \emph{Sí, con matices} (\autor{Chetty}, 2013, frente a \autor{Shiller}, 2013): formula y contrasta hipótesis, como la medicina. Exige abstracción, contraste empírico y separar positivo de normativo.
\item \emph{Frente a las naturales}: (1) experimenta menos; (2) predice peor: los sujetos reaccionan (\autor{Lucas}); (3) otro equilibrio (sendas óptimas, no punto fijo); (4) más desacuerdo $\Rightarrow$ escuelas. \emph{Frente a las sociales}: más formalización y núcleo compartido.
\item \autor{Keynes} a \autor{Harrod} (4/7/1938): «esencialmente una ciencia moral y no una ciencia natural»; \autor{Atkinson} (2009) la toma de título.
\end{itemize}

\sub{II.2. Método científico y economía\label{rep:II.2}}
\begin{itemize}
\item \concepto{Método científico}: procedimientos para investigar causas, obtener conocimiento nuevo y contrastarlo; une razón y observación.
\item \emph{Hipotético-deductivo} (\autor{Descartes}, 1637): general $\to$ particular; axiomas $\to$ conclusiones $\to$ teoremas. Rigor, pero nada nuevo sobre el mundo y premisas quizá falsas. Ej.: teoría del consumidor \vertema{3.A.8}, 1TFEB \vertema{3.A.22}.
\item \emph{Inductivo} (\autor{Bacon}, 1620): particular $\to$ general. \emph{Problema de \autor{Hume}} (1748): ningún n.º de casos garantiza el siguiente; observación nunca neutral. Pollo de \autor{Russell} (1912): «a todo cerdo le llega su San Martín». Ej.: Harvard en OI \vertema{3.A.15}; preferencia revelada (\autor{Samuelson}, 1938); RCT.
\item \autor{Kant} (1781): \emph{a priori} + \emph{a posteriori}: se discute la proporción.
\item \emph{Clásicos}: \autor{Smith} combina; \autor{Ricardo}, 1.\textsuperscript{er} constructor de \concepto{modelos}: \concepto{vicio ricardiano} (\autor{Schumpeter}): conclusiones prácticas sobre supuestos endebles; \autor{Mill} (1836): ciencia \emph{abstracta}, deductiva, \emph{a priori}, salvo causas perturbadoras: origen del \emph{homo oeconomicus} (nombre de sus críticos); \autor{Marx}: de lo abstracto a lo concreto, leyes de un modo de producción.
\item \concepto{Methodenstreit} \vertema{3.A.3}: \autor{Menger}, \emph{Untersuchungen} (1883), teoría abstracta vs. \autor{Schmoller} (reseña, 1883; economía inductiva, histórica); réplica \emph{Irrthümer} (1884). «Cada loco con su tema». Deducción vs. inducción; individualismo vs. holismo. Puente: tipos ideales de \autor{Weber} (1904). Siglo XX: \autor{Robbins} (verdades evidentes) vs. \autor{Hutchison} (1938, tautologías); \autor{Machlup} (1955): contrastar conclusiones, no cada supuesto.
\end{itemize}
\textbf{\color{tceemorado}5 respuestas del siglo XX}
\begin{itemize}
\item \emph{a) Positivismo lógico} (Círculo de Viena, años 20): \concepto{principio de verificación}; en economía, \autor{Hutchison} (1938). Límite: una ley universal no se verifica nunca.
\item \emph{b) \autor{Popper}} (1934): asimetría: ningún caso verifica, uno refuta (cisne negro). «La excepción confirma la regla»: para él, la refuta. \concepto{Falsabilidad}, criterio de demarcación: prohibir algo observable. Conjeturas y refutaciones (1963): teorías \emph{corroboradas}, no verdaderas. En economía: \autor{Friedman} (1953), afín pero instrumentalista; \autor{Blaug} (1992): se predica, no se practica. \emph{Problemas}: tesis de \concepto{\autor{Duhem}-\autor{Quine}}: la hipótesis se contrasta con auxiliares (\emph{ceteris paribus}); \autor{Popper}: sin \emph{estratagemas inmunizadoras}; tecnología; historia de la ciencia. ¿Refutó 2008 los EGDE? Pudo fallar un auxiliar (sin bancos).
\item \emph{c) \autor{Kuhn}} (1962): describe, no prescribe. \concepto{Paradigma}: logros que una comunidad reconoce como modelo de qué problemas estudiar y cómo resolverlos.
\end{itemize}
\graficorepaso[0.9\linewidth]{kuhn}{Ciencia normal (enigmas) $\to$ anomalías $\to$ crisis $\to$ revolución $\to$ nueva ciencia normal}
\begin{itemize}
\item[] \concepto{Inconmensurabilidad}: sin experimento neutral; conversión, no votación; domina no siempre la mejor. «La ciencia avanza de funeral en funeral» (\autor{Planck}). Economía: marginalista discutida (proceso de $\approx$20 años, \autor{Blaug}, 1972); keynesiana, de manual (pero síntesis); \autor{Lucas}, sí.
\item \emph{d) \autor{Lakatos}} (1970): \concepto{programa de investigación}: teorías con un mismo núcleo. \concepto{Núcleo firme}: irrefutable por decisión metodológica; \concepto{cinturón protector}: hipótesis auxiliares que se contrastan y cambian. Heurística \emph{negativa} (no tocar el núcleo: «sostenella y no enmendalla») y \emph{positiva} (cómo desarrollar el cinturón).
\end{itemize}
\graficorepaso[0.95\linewidth]{lakatos}{En gris, la nueva macro clásica: núcleo (optimización, vaciado, ER); cinturón: islas $\to$ ciclo real $\to$ EGDE con fricciones}
\begin{itemize}
\item[] \emph{Progresivo} (teórico: predice hechos nuevos; empírico: se confirman) vs. \emph{degenerativo} (parches \emph{ad hoc}). Metodología otra vez normativa. \autor{Latsis} (1976); NMC, programa progresivo \vertema{3.A.6}.
\item \emph{e) \autor{Feyerabend}} (1975): anarquismo, «todo vale» (\autor{Galileo}). \emph{Retórico}: \autor{McCloskey} (1983): se convence también con analogías y relatos. \emph{Sociológico}: \autor{Mannheim} (1929); \autor{Azoulay} \emph{et al.} (2019): muere una estrella, entran ideas de fuera. \autor{Schumpeter} (1949): \emph{visión} con ideología, que el análisis depura.
\item \emph{Balance}: falsacionistas de palabra, lakatosianos de hecho. \autor{Kuhn} describe vuelcos (\autor{Keynes}, \autor{Lucas}); \autor{Lakatos}, ajustes (NEK, RBC, EGDE tras 2008).
\end{itemize}
\begin{center}\footnotesize\setlength{\tabcolsep}{2.5pt}
\begin{tabular}{@{}>{\raggedright\arraybackslash}p{0.2\linewidth}>{\raggedright\arraybackslash}p{0.27\linewidth}>{\raggedright\arraybackslash}p{0.16\linewidth}>{\raggedright\arraybackslash}p{0.33\linewidth}@{}}
\toprule
\textbf{Escuela} & \textbf{Anomalía} & \textbf{¿\autor{Kuhn}?} & \textbf{\autor{Lakatos}}\\
\midrule
Clásicos 1776 & De dónde sale la riqueza & Fundacional & Valor trabajo, clases, acumulación\\
Marginal. 1871 & Agua y diamantes & Discutida & Núcleo nuevo: UMg, optimizador\\
\autor{Keynes} 1936 & Paro de la Depresión & Sí & Demanda efectiva, incertidumbre\\
Síntesis-monet. & Inflación con paro & No & Monetarismo: cinturón (Phillips con expectativas)\\
NMC 1972- & Estanflación & Sí (\autor{Lucas}) & Optimización, vaciado, ER; progresivo\\
NEK 1977- & Rigideces & No & Mismo núcleo; cinturón: CI, rigideces\\
Conductual 1979- & Desviaciones sistemáticas & No, aún & Amplía el cinturón\\
\bottomrule
\end{tabular}
\end{center}

\sub{II.3. Debates metodológicos\label{rep:II.3}}
\begin{itemize}
\item \emph{Parcial} (\autor{Marshall}, \emph{ceteris paribus}, ya en \autor{Molina}, 1597) vs. \emph{general} (\autor{Walras}) \vertema{3.A.21}: no compiten. «Los árboles no dejan ver el bosque». \autor{Friedman} (1949): «máquina para descubrir verdades concretas».
\item \emph{Matemáticas}: \autor{Bernoulli} (1738), \autor{Cournot} (1838), \autor{Samuelson} (1947), \autor{Arrow}-\autor{Debreu}; juegos: \autor{von Neumann}-\autor{Morgenstern} (1944), \autor{Nash} (1950), \autor{Selten} (1965), \autor{Harsanyi} (1967). A favor: rigor, generalidad. En contra: formalismo vacío (\autor{Leontief}, 1971); belleza por verdad (\autor{Krugman}, 2009); \emph{mathiness} (\autor{Romer}, 2015). \autor{Marshall}: traducir y «quemar las matemáticas» (carta a \autor{Bowley}, 1906).
\item \emph{Econometría}: contrastar, predecir, evaluar; Nobel 1969 \autor{Frisch}-\autor{Tinbergen}. \autor{Tinbergen} (1939) y crítica de \autor{Keynes}; \autor{Haavelmo} (1944), Cowles, \autor{Klein}-\autor{Goldberger} (1955); \autor{Koopmans} (1947): «medición sin teoría»; \autor{Granger} (1969): precedencia predictiva. \concepto{Crítica de \autor{Lucas}} (1976): coeficientes = reglas de decisión que cambian con la política: predicen, no evalúan. Respuestas: VAR (\autor{Sims}, 1980), calibración (\autor{Kydland}-\autor{Prescott}).
\item \emph{Supuestos}: \autor{Friedman} (1953): cuenta la predicción; billar y empresas \emph{como si}: \concepto{instrumentalismo} («gato blanco o gato negro»). \autor{Samuelson} (1963): \emph{F-twist}. \autor{Romer} (2016): supuestos cómodos pero increíbles. Modelos de 2007 sin sector financiero.
\item \emph{Modelo}: \autor{Lucas} (1980), economía artificial. \emph{7 cualidades} (\autor{Gabaix}-\autor{Laibson}, 2008): parsimonia, tratabilidad, profundidad conceptual, generalidad (las 4 aceptadas); falsabilidad, consistencia empírica, precisión predictiva. \autor{Lucas}: + parámetros estructurales.
\item \concepto{Individualismo metodológico} (marginalistas) vs. \concepto{holismo} (\autor{Marx}, escuela histórica, institucionalistas); agente representativo, \autor{Kirman} (1992). \emph{Pluralismo}: \autor{Caldwell} (1982), \autor{Rodrik} (2015): colección de modelos.
\end{itemize}

% ==========================================================================
\bloque[11,5]{III. Debates actuales; economía conductual\label{rep:III}}
\begin{ideaclave}
Todo va hacia los datos: personas con límites medidos en el laboratorio en vez del agente perfectamente racional, y una economía empírica experimental y cuasiexperimental. La conductual amplía el cinturón; 2008 obliga a ampliar los modelos macro.
\end{ideaclave}

\sub{III.1. Racionalidad y economía conductual\label{rep:III.1}}
\begin{itemize}
\item \concepto{Elección racional}: preferencias completas, transitivas y estables, interés propio, lo mejor posible \vertema{3.A.8}; versión extrema: expectativas racionales \vertema{3.A.6}. \autor{Friedman}: vale si se actúa \emph{como si}.
\item \concepto{Racionalidad limitada}: \autor{Simon} (1955): resultados \emph{satisfactorios}; Nobel 1978. \autor{Tversky}-\autor{Kahneman} (1974): heurísticas y sesgos \emph{sistemáticos} (predecibles); perspectivas (1979). \autor{Thaler} (1980).
\item \concepto{Economía conductual}: cómo deciden de hecho las personas cuando límites cognitivos, autocontrol y preferencias sociales las apartan de la racionalidad estándar; métodos de la psicología y la economía. Documenta desviaciones y las incorpora.
\item \emph{2 sistemas} (\autor{Kahneman}, 2011; nombres de \autor{Stanovich}-\autor{West}): \emph{1} rápido, automático; \emph{2} lento, deliberado.
\item \emph{Mapa} (\autor{DellaVigna}, 2009): preferencias no estándar (tiempo, riesgo, sociales); creencias (exceso de confianza, heurísticas); decisiones (marco, cuentas mentales, atención, opción por defecto).
\item \emph{$\beta$-$\delta$: idea}. Exponencial (\autor{Samuelson}, 1937): mismo descuento entre periodos consecutivos; único dinámicamente consistente (\autor{Strotz}, 1955) \vertema{3.A.29}. Evidencia (\autor{Thaler}, 1981): descuento altísimo a corto, bajo a largo: \concepto{hiperbólico}. «Una cosa es predicar y otra dar trigo».
\item \emph{Supuestos}: hogar de vida infinita, $u$ creciente y cóncava, sin incertidumbre; $\delta\in(0,1)$ exponencial, $\beta\in(0,1)$ \concepto{sesgo por el presente} (\autor{Laibson}, 1997):
\funcionclave{U_t=u(C_t)+\beta\sum_{s=1}^{\infty}\delta^{s}\,u(C_{t+s})}
\item[] Hoy-mañana: $\beta\delta$; entre 2 futuros: $\delta$. $\beta=1$: exponencial. Aproxima el \emph{hiperbólico generalizado} (\autor{Loewenstein}-\autor{Prelec}, 1992), $D(t)=(1+at)^{-\gamma/a}$.
\item \emph{Por qué no se cumple el plan}, RMS entre $t+1$ y $t+2$:
\begin{formula}
\centering desde $t$: $\dfrac{u'(C_{t+1})}{\delta\,u'(C_{t+2})}$;\quad desde $t+1$: $\dfrac{u'(C_{t+1})}{\beta\delta\,u'(C_{t+2})}$
\end{formula}
\item[] $\beta<1$: al llegar $t+1$, más consumo y menos ahorro de lo planeado.
\item \emph{Ejemplo} ($\beta=0{,}7$, $\delta=0{,}99$, semanas): hoy, 100 vs. 110 en 1 semana: $100>0{,}7\cdot0{,}99\cdot110\approx76{,}2$: 100. Semana 52 vs. 53: $100<0{,}99\cdot110\approx108{,}9$: esperar. Llega la 52: 100. Sin información nueva.
\end{itemize}
\graficorepaso[0.85\linewidth]{descuento}{Salto en $t=1$ (el $\beta$), luego exponencial. 1 € a 1 año: 0,94 exp. vs. 0,67 cuasihip.; a 30 años: 0,15 vs. 0,19 (\autor{Angeletos} \emph{et al.}, 2001)}
\begin{itemize}
\item \emph{Implicaciones}: \concepto{inconsistencia temporal} (juego entre «yoes»); infraahorro y procrastinación; \emph{demanda de compromiso}: \emph{sofisticado} se ata (\autor{Ulises}); \emph{ingenuo} procrastina más (\autor{O'Donoghue}-\autor{Rabin}, 1999).
\item \emph{Extensiones}: consumo sensible a la renta corriente \vertema{3.A.33}; búsqueda de empleo \vertema{3.A.26}; política: \autor{Kydland}-\autor{Prescott} (1977) \vertema{3.A.36}.
\item \emph{Evidencia}: descuento cae con el horizonte (\autor{Frederick} \emph{et al.}, 2002); gimnasio: $>$70 \$/mes, 4,3 visitas, $>$17 \$/visita vs. bono a 10 (\autor{DellaVigna}-\autor{Malmendier}, 2006); Filipinas: cuenta bloqueada, saldo +81 p.p. (\autor{Ashraf}-\autor{Karlan}-\autor{Yin}, 2006); Kenia: fertilizante tras la cosecha (\autor{Duflo}-\autor{Kremer}-\autor{Robinson}, 2011).
\item \emph{Valoración}: +1 parámetro, explica el compromiso, conserva la estructura. Pero $\beta$ difícil de medir (fungibilidad: con esfuerzo sí, con dinero menos; \autor{Andreoni}-\autor{Sprenger}, 2012); se confunde con liquidez; ¿ingenuo o sofisticado? ¿Qué «yo» cuenta para el bienestar?
\item \emph{Riesgo: perspectivas} (\autor{Kahneman}-\autor{Tversky}, 1979) \vertema{3.A.10}: punto de referencia $r$; \concepto{aversión a las pérdidas} ($\lambda\approx2{,}25$, 1992); sensibilidad decreciente (forma de S); probabilidades no lineales.
\end{itemize}
\graficorepaso[0.7\linewidth]{funcion-valor}{Cóncava en ganancias, convexa en pérdidas, más pendiente en pérdidas: $v(-a)=-\lambda v(a)$}
\begin{itemize}
\item[] Explica: \emph{efecto dotación} (tazas, $\approx$2×; \autor{Kahneman}-\autor{Knetsch}-\autor{Thaler}, 1990); \emph{statu quo} (\autor{Samuelson}-\autor{Zeckhauser}, 1988; «virgencita, que me quede como estoy»); rigidez nominal a la baja (1986); felicidad y renta (\autor{Easterlin}, 1974; \autor{Stevenson}-\autor{Wolfers}, 2008).
\item \emph{Preferencias sociales}: \autor{Smith} (1759); aversión a la desigualdad (\autor{Fehr}-\autor{Schmidt}, 1999); ERC (\autor{Bolton}-\autor{Ockenfels}, 2000).
\item \emph{Heurísticas} (1974): \emph{representatividad}; \emph{disponibilidad} («el que se quema con leche\ldots»); \emph{anclaje} (ruleta 10 o 65 $\to$ 25 o 45~\%). Exceso de confianza; \concepto{sesgo retrospectivo} (\autor{Fischhoff}, 1975): ¡el refrán del enganche!
\item \emph{Decisiones}: \concepto{efecto marco} (epidemia de 600: 72~\% seguro si «salvar», 78~\% riesgo si «morir»; 1981); \concepto{cuentas mentales} (\autor{Thaler}, 1985): teatro, 88~\% vs. 46~\%.
\item \concepto{Paternalismo libertario} (\autor{Thaler}-\autor{Sunstein}, 2003) y \concepto{empujón} (\emph{nudge}, 2008): cambiar la arquitectura de la elección sin prohibir \vertema{3.A.20}; no hay arquitectura neutral. Ej.: inscripción automática (\autor{Madrian}-\autor{Shea}, 2001; R.\,Unido); órganos (\autor{Johnson}-\autor{Goldstein}, 2003); \emph{Save More Tomorrow}: 3,5 $\to$ 13,6~\% en 40 meses (2004). \emph{Behavioural Insights Team}. Finanzas conductuales: límites al arbitraje (\autor{Shleifer}-\autor{Vishny}, 1997) \vertema{3.B.23}.
\item \emph{Valoración}: + realismo, políticas baratas; caja de herramientas (\autor{Gabaix}, 2020). Críticas: lista, no teoría (\autor{Fudenberg}, 2006); laboratorio $\neq$ mercado (\autor{List}, 2003; \autor{Levitt}-\autor{List}, 2007); replicación (psicología 36~\%; economía 11/18 = 61~\%); \emph{nudges} de gobiernos 1,4 p. vs. 8,7 publicados (\autor{DellaVigna}-\autor{Linos}, 2022); racionalidad ecológica (\autor{Gigerenzer}); paternalismo.
\end{itemize}

\sub{III.2. Experimentos en economía\label{rep:III.2}}
\begin{itemize}
\item \emph{Laboratorio}: V.~\autor{Smith} (1962): convergencia al equilibrio competitivo; Nobel 2002. \emph{Ultimátum} (\autor{Güth} \emph{et al.}, 1982): \autor{Nash} PS = mínimo \vertema{3.A.14}; oferta media 40~\%, rechazo 16~\%. \emph{Dictador} (\autor{Forsythe} \emph{et al.}, 1994): da 28~\% (\autor{Engel}, 2011): parte de lo del ultimátum es miedo al rechazo.
\item \emph{Naturales}: \autor{Card}-\autor{Krueger} (1994): NJ, 4,25 $\to$ 5,05 \$ (+19~\%) vs. Pensilvania: el empleo no cae \vertema{3.A.25}. Nobel 2021.
\item \emph{RCT}: $Y_i(1)$, $Y_i(0)$; solo se ve uno; comparar mezcla efecto y \concepto{sesgo de selección}; el sorteo lo elimina (\autor{Angrist}-\autor{Pischke}, 2009) \vertema{3.B.28}. «Ver para creer». Nobel 2019. Validez interna sí, externa no.
\item Modelos \emph{estructurales} (teoría $\to$ datos) vs. \emph{experimentales} (datos $\to$ teoría).
\end{itemize}

\sub{III.3. Innovaciones metodológicas\label{rep:III.3}}
\begin{itemize}
\item \emph{Giro empírico}: teoría pura 51~\% (1963) $\to$ 19~\% (2011); empírico con datos propios 9 $\to$ 34~\% (\autor{Hamermesh}, 2013). \concepto{Revolución de la credibilidad} (\autor{Angrist}-\autor{Pischke}, 2010, tras \autor{Leamer}, 1983): diseñar para identificar un efecto causal (DiD, VI, regresión discontinua).
\item \emph{¿Evolución o revolución tras 2008?} Revolución: \autor{Romer} (2016), \autor{Stiglitz} (2018), modelos de agentes (\autor{Leijonhufvud}, 2006). Intermedio: \autor{Farmer} (2010), equilibrios múltiples. Evolución: \autor{Blanchard} (2018), \autor{Reis} (2018), \autor{Galí} (2018), \emph{Rebuilding Macroeconomic Theory} (2018). Lectura de \autor{Lakatos}: núcleo intacto, cinturón más ancho; sin paradigma rival.
\item \emph{Heterogeneidad}: HANK (\autor{Kaplan}-\autor{Moll}-\autor{Violante}, 2018) \vertema{3.A.7}; microdatos (\autor{Piketty}-\autor{Saez}, 2003; \autor{Chetty} \emph{et al.}, 2014); granularidad (\autor{Gabaix}, 2011).
\item \emph{Complejidad}: Santa Fe (\autor{Arthur}, 1999); modelos basados en agentes; caos (\autor{Grandmont}, 1985); redes (\autor{Acemoglu} \emph{et al.}, 2012). Choca con parsimonia y falsabilidad.
\item \emph{Aprendizaje automático}: bueno para $\hat y$, no para $\hat\beta$ (\autor{Mullainathan}-\autor{Spiess}, 2017; \autor{Athey}-\autor{Imbens}, 2019); \autor{Hume} y \autor{Lucas} siguen ahí. \emph{Replicación}: potencia baja (\autor{Ioannidis} \emph{et al.}, 2017); macro: 1/3 se reproduce (\autor{Chang}-\autor{Li}, 2022); registro previo, datos y código.
\end{itemize}

% ==========================================================================
\bloque[1]{Conclusión}
\begin{itemize}
\item \textbf{Recapitulación}: elegir con recursos escasos (de la riqueza a cualquier elección); método hipotético-deductivo contrastado, mejor con \autor{Kuhn} y \autor{Lakatos} que con falsacionismo estricto; hoy a prueba con psicología, experimentos y microdatos; $\beta$-$\delta$ como modelo.
\item \textbf{Relevancia}: el método decide qué políticas se evalúan; por \autor{Lucas}, modelos estructurales; por la conductual, inscripción automática; por los experimentos, qué se financia.
\item \textbf{Extensiones}: escuelas \vertema{3.A.2}--\vertema{3.A.7}; riesgo \vertema{3.A.10}; juegos \vertema{3.A.14}; intertemporal \vertema{3.A.29}; consumo \vertema{3.A.33}; finanzas \vertema{3.B.23}; evaluación \vertema{3.B.28} \vertema{4.B.6}.
\item \textbf{Opinión}: contra la «visión acumulativa», las ideas vuelven (\autor{Simon}, la inducción, equilibrios múltiples; \autor{Roncaglia}, 2006): pluralismo de modelos (\autor{Rodrik}).
\item \textbf{Idea final}: a toro pasado todos la vemos venir; lo científico es convertir cada anomalía en una pregunta contrastable. Respuesta a la reina: no un paradigma nuevo, sino un cinturón protector más ancho y realista.
\end{itemize}

% ==========================================================================
\bloque{Autores y fechas}
\hito{1615}{\autor{Montchrestien}: «economía política»}
\hito{1620}{\autor{Bacon}: \emph{Novum Organum}}
\hito{1637}{\autor{Descartes}: \emph{Discurso del método}}
\hito{1748}{\autor{Hume}: problema de la inducción}
\hito{1776}{\autor{Smith}: \emph{La riqueza de las naciones}}
\hito{1781}{\autor{Kant}: \emph{Crítica de la razón pura}}
\hito{1836}{\autor{Mill}: definición; 1843, \emph{Lógica}}
\hito{1883}{\emph{Methodenstreit}: \autor{Menger}-\autor{Schmoller}}
\hito{1890}{\autor{Marshall}: \emph{Principles}}
\hito{1891}{J.~N.~\autor{Keynes}: positiva, normativa, arte}
\hito{1932}{\autor{Robbins}: fines y medios escasos}
\hito{1934}{\autor{Popper}: falsabilidad}
\hito{1938}{\autor{Hutchison}; \autor{Keynes}: «ciencia moral»}
\hito{1944}{\autor{Haavelmo}; \autor{von Neumann}-\autor{Morgenstern}}
\hito{1949}{\autor{Mises}: praxeología}
\hito{1953}{\autor{Friedman}: \emph{Methodology}; \autor{Myrdal}}
\hito{1955}{\autor{Simon}: racionalidad limitada; \autor{Strotz}}
\hito{1962}{\autor{Kuhn}: paradigmas; V.~\autor{Smith}}
\hito{1970}{\autor{Lakatos}: programas (recop. 1978)}
\hito{1974}{\autor{Tversky}-\autor{Kahneman}: heurísticas}
\hito{1975}{\autor{Feyerabend}: «todo vale»}
\hito{1976}{\autor{Lucas}: crítica; \autor{Becker}}
\hito{1979}{\autor{Kahneman}-\autor{Tversky}: perspectivas}
\hito{1980}{\autor{Thaler}; \autor{Sims}: VAR}
\hito{1982}{\autor{Güth} \emph{et al.}: ultimátum}
\hito{1983}{\autor{McCloskey}: retórica; \autor{Leamer}}
\hito{1992}{\autor{Loewenstein}-\autor{Prelec}: hiperbólico}
\hito{1994}{\autor{Card}-\autor{Krueger}: salario mínimo}
\hito{1997}{\autor{Laibson}: $\beta$-$\delta$}
\hito{2003}{\autor{Thaler}-\autor{Sunstein}: paternalismo libertario}
\hito{2008}{\emph{Nudge}; \autor{Gabaix}-\autor{Laibson}: 7 cualidades}
\hito{2010}{\autor{Angrist}-\autor{Pischke}: credibilidad}
\hito{2016}{\autor{Romer}: \emph{The Trouble with Macro}}
\textbf{Nobel}: 1969 \autor{Frisch}-\autor{Tinbergen}; 1974 \autor{Myrdal}-\autor{Hayek}; 1978 \autor{Simon}; 1992 \autor{Becker}; 1994 \autor{Nash}-\autor{Selten}-\autor{Harsanyi}; 2002 \autor{Kahneman}-V.~\autor{Smith}; 2017 \autor{Thaler}; 2019 \autor{Banerjee}-\autor{Duflo}-\autor{Kremer}; 2021 \autor{Card}-\autor{Angrist}-\autor{Imbens}.

% ==========================================================================
\bloque{Datos clave}
\begin{itemize}
\item Pensiones con inscripción automática, R.\,Unido: asalariados con derecho que ahorran en un plan de empresa, \dato{55 \% (2012) $\to$ 90 \% (2025)}{DWP, 2026}.
\item Teoría pura en las 3 grandes revistas \dato{51 \% $\to$ 19 \%}{Hamermesh, 2013; 1963-2011}; empírico con datos propios \dato{9 $\to$ 34 \%}{ídem}.
\item Replicación: \dato{36 \%}{OSC, 2015; psicología}; \dato{11 de 18 (61 \%), 2/3 del efecto}{Camerer \emph{et al.}, 2016; economía}.
\item Prima de las acciones EE.\,UU. 1928-2025: \dato{$\approx$6,7 p. sobre letras; 5,5 sobre bonos a 10 años}{Damodaran, 2026}.
\end{itemize}

% ==========================================================================
\bloque{Preguntas del tribunal}
\begin{enumerate}
\item \textbf{Defina «paradigma».} Logros que una comunidad reconoce durante un tiempo como modelo de qué problemas estudiar y cómo (\autor{Kuhn}, 1962); domina en la ciencia normal hasta la crisis.
\item \textbf{¿Verificabilidad o falsabilidad?} Verificar: positivismo lógico; una ley universal no se verifica. \autor{Popper}: un caso refuta. Crítica: \autor{Duhem}-\autor{Quine} (auxiliares), historia de la ciencia (\autor{Kuhn}).
\item \textbf{¿Revoluciones científicas en economía?} \autor{Keynes} de manual, pero integrado en la síntesis; \autor{Lucas}, sí; marginalista, discutida; NEK y conductual: cinturón (\autor{Lakatos}).
\item \textbf{Un mercado desde la positiva y la normativa.} Salario mínimo: ¿efecto en el empleo? (positiva) / ¿es justo? (normativa) / ¿qué nivel? (arte).
\item \textbf{\autor{Friedman} y los supuestos; \autor{Romer} (2016).} Instrumentalismo: cuenta la predicción (\emph{como si}); \autor{Samuelson}: \emph{F-twist}; \autor{Romer}: supuestos increíbles que blindan los modelos.
\item \textbf{Crítica de \autor{Lucas}.} Los modelos macroeconométricos de la síntesis (Cowles, \autor{Klein}-\autor{Goldberger}) recogen reglas que cambian con la política: no sirven para evaluarla $\Rightarrow$ microfundamentar.
\item \textbf{¿Qué es el imperialismo económico?} Llevar el método (optimización, equilibrio) fuera del mercado: \autor{Becker}; \autor{Lazear} (2000). Contrario: la conductual importa de la psicología.
\end{enumerate}

\begin{trampa}
\begin{itemize}
\item \autor{Robbins}: lo alternativo son los \emph{usos de los medios}; la definición es suya, no de \autor{Samuelson}.
\item «Ciencia moral»: \autor{Keynes} a \autor{Harrod} (1938), no \autor{Atkinson}. «Lo que hacen los economistas»: atribuida a \autor{Viner}, publicada por \autor{Boulding} (1941).
\item \emph{Homo oeconomicus}: abstracción de \autor{Mill}, nombre de sus críticos.
\item Praxeología: \autor{Mises} (1949), no \autor{Menger}; \emph{Methodenstreit}: \autor{Menger} vs. \autor{Schmoller}, no \autor{Walras}.
\item \autor{Friedman}: instrumentalista; \autor{Popper}: no. En \autor{Lakatos} se contrasta el \emph{cinturón}, no el núcleo. \emph{Proofs and Refutations} (1976): filosofía de las matemáticas.
\item Hiperbólico generalizado: \autor{Loewenstein}-\autor{Prelec} (1992); cuasihiperbólico: \autor{Laibson} (1997); inconsistencia: \autor{Strotz} (1955). En $\beta$-$\delta$, $\beta$ = sesgo por el presente (en el resto del temario, $\beta$ = factor de descuento).
\item Ultimátum: \autor{Güth}, \autor{Schmittberger} y \autor{Schwarze} (1982), no \autor{Thaler}. \autor{Tversky} sin Nobel (murió en 1996).
\item Causalidad de \autor{Granger}: precedencia predictiva, no causalidad estricta.
\end{itemize}
\end{trampa}

\end{columnas}
\end{document}
